\documentclass[11pt]{tg_usta}
\labelpaper[year =2026 , month =6]
\usepackage{anysize}
\marginsize{1in}{1in}{1in}{1in}
\usepackage{amsmath,amsthm,amssymb}
\usepackage{graphicx}
\usepackage{booktabs}
\usepackage{array}
\usepackage{url}
\decimalpoint

\begin{document}
\title[maintitle = Bitácora de uso de Inteligencia Artificial,
       secondtitle = Proyecto Centinela --- Fase 1: línea base con un MLP,
       shorttitle = Bitácora de IA --- Centinela Fase 1]

\begin{authors}
\author[firstname = Estefanny,
        surname = Ruíz González,
        affiliation = {Estudiante},
        email = estefannyruiz@usantotomas.edu.co]
\author[firstname = Miguel,
        surname = Alarcón Ojeda,
        affiliation = {Estudiante},
        email = miguel.alarcon@usantotomas.edu.co]
\end{authors}

\section{Propósito y forma de uso}

Esta bitácora documenta el uso de una herramienta de inteligencia artificial conversacional (Claude, de Anthropic) como apoyo durante la Fase 1 del Proyecto Centinela. El uso siguió el criterio de \emph{preguntar, verificar, corregir y rendir cuentas}: la IA se empleó para explicar conceptos, proponer decisiones de diseño y generar borradores de código y documentos, pero cada sugerencia fue evaluada por los autores antes de incorporarse. Las decisiones se clasifican en \textbf{aceptadas} (incorporadas tal cual), \textbf{modificadas} (incorporadas con cambios) y \textbf{rechazadas} (descartadas), con su justificación. Ninguna salida de la IA se adoptó sin revisión.

\section{Sugerencias aceptadas}

\begin{itemize}
\item \textbf{Reformulación del problema a clasificación binaria.} La IA propuso convertir el problema descriptivo del proyecto de grado (medir el poder adquisitivo) en la predicción de un evento de riesgo binario (``el salario real está cayendo''), adecuado para el MLP con salida sigmoide. Se aceptó por ser coherente con la consigna del módulo y con la idea de alerta temprana.
\item \textbf{Exclusión de variables por fuga de datos.} La IA identificó que \texttt{smlv\_real}, \texttt{smlv\_real\_var\_anual} y \texttt{poder\_adq\_index} no podían usarse como variables de entrada por ser la base de la etiqueta (la última, un reescalado de \texttt{smlv\_real}). Se aceptó tras verificar la equivalencia numérica en los datos.
\item \textbf{Partición cronológica.} La IA recomendó dividir train/val/test en orden temporal (no aleatorio) por tratarse de una serie de tiempo. Se aceptó porque evita la fuga temporal, y se mantuvo la proporción 70/15/15 indicada por el docente.
\item \textbf{Uso de \texttt{BCEWithLogitsLoss} y \texttt{pos\_weight}.} Se aceptó la salida en logit (sin sigmoide en el modelo) combinada con esta pérdida por estabilidad numérica, y el ponderado de la clase minoritaria para el desbalance ($\sim$9.6\,\%).
\item \textbf{Framework PyTorch.} Se aceptó la recomendación de usar PyTorch por alinearse con el código de referencia del curso.
\end{itemize}

\section{Sugerencias modificadas}

\begin{itemize}
\item \textbf{Corrección del porcentaje de desbalance.} La IA escribió en una sección del notebook que la clase de riesgo era el 9.6\,\%, cuando en ese punto (antes de la limpieza de datos) correspondía calcularlo sobre el total de 324 registros, es decir 9.3\,\%. Los autores detectaron y corrigieron la imprecisión, dejando 9.3\,\% antes de la limpieza y 9.6\,\% después, según la población medida en cada etapa.
\item \textbf{Definición de la segunda modalidad.} La IA sugirió inicialmente una modalidad de imágenes (luminosidad nocturna satelital) para la fusión de la Fase 2. Se modificó: dado que el proyecto se centra en datos económicos, se declaró la modalidad secuencial (redes recurrentes sobre la misma serie) como la segunda modalidad, dejando la consulta de este punto al docente.
\item \textbf{Ubicación de las extensiones propuestas por la asesora.} La IA consideró incluir el índice de Brecha y el enfoque DTW en la propuesta de datos. Se modificó su ubicación: se trasladaron a la sección de ``trabajo futuro'' del informe técnico-ético, por estar fuera del alcance de esta fase.
\item \textbf{Selección final del conjunto de variables.} Sobre varias features candidatas propuestas por la IA, los autores acotaron el conjunto a las de dinámica del IPC y la brecha nominal-inflación, tras revisar las correlaciones, descartando las de baja señal.
\end{itemize}

\section{Sugerencias rechazadas}

\begin{itemize}
\item \textbf{Agrupamiento DTW multi-ciudad en esta fase.} La IA planteó la posibilidad de un agrupamiento jerárquico por \emph{Dynamic Time Warping}. Se rechazó para la Fase 1 por exceder su alcance (que exige una sola ciudad y un MLP supervisado), conservándose únicamente como trabajo futuro.
\item \textbf{Índice de Brecha (canasta/SMLV) como variable inmediata.} Se rechazó incorporarlo ahora porque la base actual no contiene el costo de la canasta familiar; requeriría consolidar otra fuente, propio de una fase posterior.
\item \textbf{Variable TRM como predictor.} Se rechazó por mostrar correlación casi nula con la etiqueta, para no añadir ruido a un modelo con muestra reducida.
\item \textbf{Diferenciación de dos enfoques individuales.} Se había planteado un enfoque distinto por integrante; se descartó al confirmarse que la entrega es grupal, unificando un solo conjunto de entregables.
\end{itemize}

\section{Reflexión sobre el uso crítico}

El principal aporte de la herramienta fue acelerar la comprensión de conceptos (sobreajuste, fuga de datos, partición temporal) y la generación de borradores, pero las decisiones de fondo --qué variables usar, cómo evaluar ante el desbalance, qué priorizar éticamente-- se tomaron y verificaron de forma independiente. El episodio de la corrección del porcentaje ilustra la necesidad de no aceptar las salidas sin contraste: una cifra plausible puede ser imprecisa según el contexto. Los autores asumen la responsabilidad sobre el contenido final de todos los entregables y están en capacidad de sustentar cada decisión sin asistencia.

\end{document}
